% ============================================================================
%  版式设置：preamble/layout.tex
%  职责：页面几何、行距、缩进、标题与图表题注格式、浮动体比例
% ============================================================================

% --- 页面几何 ---
\geometry{top=2.5cm, bottom=2.5cm, left=2.8cm, right=2.8cm, headsep=1.0cm}

% --- 段落 ---
\setlength{\parindent}{2em}          % 中文首行缩进两个字符
\setlength{\parskip}{0pt}            % 段落之间不额外加间距（期刊标准做法）
\renewcommand{\arraystretch}{1.2}    % 表格行高，避免内容挤在一起

% --- 浮动体控制（防止图表大量漂移到文末）---
\renewcommand{\topfraction}{0.85}
\renewcommand{\bottomfraction}{0.75}
\renewcommand{\textfraction}{0.10}
\renewcommand{\floatpagefraction}{0.80}
\setcounter{topnumber}{3}
\setcounter{bottomnumber}{2}
\setcounter{totalnumber}{5}

% --- 题注格式 ---
\captionsetup{
  font        = small,
  labelfont   = bf,
  labelsep    = quad,
  figurename  = 图,
  tablename   = 表,
  justification = centering,
  width       = 0.92\linewidth
}

% --- 定理类环境 ---
\newtheoremstyle{plaincn}
  {1.0em}{1.0em}{\itshape}{}{\bfseries}{}{0.6em}{}
\theoremstyle{plaincn}
\newtheorem{theorem}{定理}[section]
\newtheorem{lemma}[theorem]{引理}
\newtheorem{corollary}[theorem]{推论}
\newtheorem{proposition}[theorem]{命题}
\newtheorem{definition}[theorem]{定义}
\newtheorem{assumption}[theorem]{假设}
\newtheorem{remark}[theorem]{注}
